% Chapter 3: Ontology Design
\section{Overview}

This chapter details the design and development of the CS education ontology that forms the core of the framework.

\section{Ontology Engineering Methodology}

\subsection{Requirements Gathering}

The ontology development process began with gathering requirements from educators, students, and domain experts. These requirements informed the conceptualization phase.

\subsection{Conceptualization}

Key concepts in CS education include:
\begin{itemize}
    \item Learning objectives and outcomes
    \item Core CS concepts and topics
    \item Pedagogical strategies
    \item Assessment mechanisms
    \item Student profiles and preferences
\end{itemize}

\section{Core Ontology Components}

\subsection{Topic Hierarchy}

The ontology defines a hierarchical organization of CS topics, from foundational concepts to advanced specializations.

\subsection{Learning Relationships}

Relationships such as \textit{prerequisite}, \textit{reinforces}, and \textit{applies-to} connect concepts and enable reasoning about learning paths.

\subsection{Metadata and Properties}

Each concept includes metadata such as difficulty level, learning time estimate, and recommended instructional strategies.

\section{Formalization}

The ontology is formally represented using OWL (Web Ontology Language), enabling automated reasoning and semantic querying.

\section{Validation}

The ontology underwent validation through expert review and comparison with existing educational standards and frameworks.
